\documentclass[sigconf,nonacm]{acmart}
\settopmatter{printfolios=true}
\usepackage{booktabs}
\usepackage{hyperref}
\newcommand{\todo}[1]{\textcolor{red}{[#1]}}

\begin{document}
\title{Artifact Appendix: AdLens: Efficient Detection of Deceptive Software Ads}
\author{Ritik Roongta, Marwan Adnan Darwish, Masoud Poorghaffar Aghdam, Rachel Greenstadt, Gunes Acar}
\maketitle

\section{Artifact Appendix}

\subsection{Abstract}
AdLens detects deceptive software ads (scareware, deceptive claims, misleading ad design) with a pipeline of
OCR (PaddleOCR), translation (\texttt{translategemma:4b}), semantic-similarity ranking, a two-model VLM/LLM
classifier ensemble (\texttt{qwen3.5:9b}, \texttt{gemma3:12b}), and a judge (\texttt{gemma4:26b}) for
classifier disagreements. The artifact contains the crawler, OCR, translation, ranking and detection code,
the SQL used for the measurements, the labeled golden datasets, a sample dataset, the semantic search tool,
Docker images, and scripts that recompute Table~2 from freshly produced pipeline output.

\subsection{Artifact check-list (meta-information)}
{\small
\begin{itemize}
  \item \textbf{Algorithm:} OCR $\rightarrow$ translation $\rightarrow$ embedding-based ranking $\rightarrow$ two-model classification $\rightarrow$ judge on disagreement.
  \item \textbf{Data set:} Labeled golden set (1{,}283 annotations, 1{,}271 images) and held-out test set (253 annotations, 251 images), both in \texttt{datasets/}; unlabeled sample in \texttt{sample\_data/}.
  \item \textbf{Run-time environment:} Linux, Python 3.12, Docker (CPU and GPU images) or a local virtualenv; Ollama \textbf{0.40.6}.
  \item \textbf{Hardware:} GPU with $\geq$24\,GB VRAM recommended for \texttt{gemma4:26b} (we used an NVIDIA L40S).
  \item \textbf{Metrics:} Per-class precision, recall, F1 and macro F1.
  \item \textbf{Output:} JSON results under \texttt{detection/results/}; Table~2 printed by \texttt{scripts/eval\_pipeline\_results.py}.
  \item \textbf{Experiments:} Table~2(a), 2(b), 2(c).
  \item \textbf{How much disk space is required?:} $\sim$40\,GB for the four model builds plus images.
  \item \textbf{How much time is needed to prepare the workflow?:} $\sim$1 hour (mostly model downloads).
  \item \textbf{How much time is needed to complete experiments?:} Held-out set $\sim$1\,h; full golden set $\sim$5.5\,h on one L40S.
  \item \textbf{Publicly available?:} Yes. Zenodo: \todo{DOI}; GitHub: \url{https://github.com/Racro/ccs-2026-AdLens}.
  \item \textbf{Code licenses:} MIT (see \texttt{LICENSE}).
\end{itemize}}

\subsection{Description}

\subsubsection{How to access}
Zenodo archive \todo{DOI} (permanent) and \url{https://github.com/Racro/ccs-2026-AdLens}. Docker images: \texttt{racro/adlens} (CPU and GPU tags, see \texttt{README.md}).

\subsubsection{Hardware dependencies}
CPU-only execution works for OCR and tiny samples but is impractically slow for the 26B judge. A single
24--48\,GB GPU is sufficient for all experiments; models are run sequentially through Ollama.

\subsubsection{Software dependencies}
Python 3.12, \texttt{requirements.txt} and \texttt{ocr/requirements.txt}, PaddlePaddle/PaddleOCR 3.x, and
\textbf{Ollama 0.40.6}. Older Ollama builds (e.g.\ 0.21) cannot load \texttt{gemma4}.

\subsubsection{Data sets}
\texttt{datasets/golden\_main\_data} (Table~2(a)/(b)) and \texttt{datasets/golden\_test\_data} (Table~2(c)),
each with \texttt{metadata.json} (violation type, \texttt{tp}/\texttt{tn} label, OCR text) and the ad images.
Twelve (resp.\ two) ads carry two labels and share one image file. Four held-out images also appear in the
main set; excluding them changes held-out macro F1 by $<$0.001 (0.9687 vs.\ 0.9683).
The 188{,}834-creative crawl itself is not redistributed (Section~\ref{sec:limits}).

\subsubsection{Models}
Model tags on the Ollama registry are mutable and results depend on the exact build \emph{and} runtime.
The recorded builds are listed with manifest SHA-256 digests in \texttt{run.md}. They can be restored with a
single command (Section~\ref{sec:setup}).

\subsection{Installation}\label{sec:setup}
\begin{enumerate}
  \item Clone the repository (or unpack the Zenodo archive) and either pull the Docker image or create a virtualenv
        (\texttt{pip install -r requirements.txt -r ocr/requirements.txt}).
  \item Install Ollama \textbf{0.40.6}, start \texttt{ollama serve} with \texttt{OLLAMA\_MODELS} pointing to an empty directory.
  \item Restore the recorded model builds (downloads the layer blobs by SHA-256, verifies them and installs the recorded manifests):
\begin{verbatim}
export OLLAMA_MODELS=$HOME/ollama_models
python scripts/restore_ollama_models.py
ollama list
\end{verbatim}
        Do \emph{not} use \texttt{ollama pull} for reproduction: it fetches the current tag, which has changed since our runs.
\end{enumerate}

\subsection{Experiment workflow}
The held-out and golden evaluations run the full pipeline with \texttt{--eval}, which restricts each
violation type to the ads labeled for it, and then score the output.
See the ``Reproducing Table~2'' section of \texttt{README.md} for the exact commands and flags.

\subsection{Evaluation and expected results}
\begin{itemize}
  \item \textbf{E1 -- Table~2(a), classifiers (golden set).} Macro F1: Qwen3.5:9b 0.929, Gemma3:12b 0.847.
  \item \textbf{E2 -- Table~2(b), judge and full cascade (golden set).} Judge macro F1 0.894 and full-cascade macro F1 0.946 (DC 0.909, SW 0.987, MD 0.940).
  \item \textbf{E3 -- Table~2(c), held-out set (253 ads).} Full-cascade macro F1 0.969 (DC 0.949, SW 0.957, MD 1.000; precision 1.000 in every class).
\end{itemize}
Run-to-run variation from sampling is visible: repeated held-out runs gave macro F1 between 0.965 and 0.969,
and judge outputs vary slightly between runs (e.g.\ 127 vs.\ 128 judged cases).
Remaining error is dominated by classifier recall: positives that both classifiers label safe never reach the judge.
Disabling the thinking trace for \texttt{gemma4} (the default in the released code) and enabling it changed the golden
full-cascade macro F1 only from 0.950 to 0.946.

All numbers above come from end-to-end runs of the released code.

\subsection{Notes and limitations}\label{sec:limits}
\begin{itemize}
  \item \textbf{Large-scale measurements (Figures~3--5, Section~5).} These are computed over the full 188{,}834-creative crawl
        and the ranked corpus derived from it. The artifact provides the code and the SQL queries (with their ID-list cardinalities), but not the
        full crawl or all ad images, so these figures are not regenerated by the Table~2 workflow.
  \item \textbf{Latency.} On our hardware the judge takes 30--60\,s per call with the Ollama default (thinking on); the paper's latency figures (Appendix~B) were measured separately and are not recomputed by the Table~2 workflow.
  \item \textbf{Configuration sensitivity.} Pulling current tags or using a newer Ollama can materially change held-out results; use the restore script.
  \item \textbf{Dataset overlap.} Four held-out images are pixel-identical to main-set images (see above); they were left in the released set and the effect is reported.
\end{itemize}

\subsection{Methodology}
Submission, reviewing and badging methodology:
\begin{itemize}
  \item \url{https://www.acm.org/publications/policies/artifact-review-and-badging-current}
  \item \url{https://sysartifacts.github.io/evaluator-guide.html}
\end{itemize}

\end{document}
